\documentclass[../main.tex]{subfiles}
\begin{document}
\chapter*{ABSTRACT}
\thispagestyle{plain}

In most universities today, nearly every student already uses Zalo. Deploying an attendance
application directly on the Zalo Mini App platform therefore brings significant convenience:
students can check in immediately without installing any additional software. Meanwhile,
traditional methods such as oral roll-call or paper sign-in sheets are time-consuming for large
classes and cannot reliably guarantee that students are genuinely present.

This thesis presents the design and implementation of a smart attendance system named
\textit{Zimo Check-in}, built on the Zalo Mini App platform. The check-in procedure is designed
in four steps: scanning a dynamic QR code that is signed with HMAC-SHA256 and refreshed every
30 seconds, verifying the student's face directly on the device, exchanging peer-to-peer
confirmation codes among nearby students, and aggregating these results into a trust policy.
Combining such complementary verification layers helps ensure accurate and honest attendance
records without requiring expensive dedicated hardware.

In terms of architecture, the product comprises a Mini App for students and lecturers, a
web-based administration dashboard for the academic office, and server-side functions running on
Firebase. It further provides supporting features such as a projector gateway that displays a
large QR code for crowded classrooms, real-time attendance monitoring, and leave-request
management. Experimental results show that the system operates stably and satisfactorily meets
the functional and security requirements defined for it.

\vspace{1em}
\noindent\textbf{Keywords:} \textit{smart attendance, Zalo Mini App, rotating QR, HMAC-SHA256, face verification, Firebase.}
\end{document}
